% Frozen formal layer — papersmith P1 (graph render). Bodies are hash-pinned; do not edit.
\begin{assumptionv}{ass:iid}[Independent sampling]
The observed data units
\[
O_1,\ldots,O_n
\]
are independent and identically distributed draws from \(P\).
\end{assumptionv}

\begin{assumptionv}{ass:bounded-outcome}[Bounded potential outcomes]
For each treatment value \(a\in\{0,1\}\),
\[
Y(a)\in[-1,1].
\]
\end{assumptionv}

\begin{assumptionv}{ass:positivity}[Positivity]
The propensity score satisfies
\[
0<e_P(X)<1
\]
\(P_X\)-almost surely.
\end{assumptionv}

\begin{assumptionv}{ass:margin}[Margin condition]
For every \(u\) with \(0<u\le u_0\),
\[
P\{0<|\tau_P(X)|\le u\}\le C_m u^\alpha .
\]
\end{assumptionv}

\begin{assumptionv}{ass:zero-effect}[Zero-effect agreement]
Either
\[
P_X\{\tau_P(X)=0\}=0,
\]
or every policy \(\pi\in\Pi\) agrees with the oracle policy \(\pi^\star_P\) \(P_X\)-almost everywhere on the set \(\{\tau_P(X)=0\}\).
\end{assumptionv}

\begin{assumptionv}{ass:policy-class}[Finite-VC policy class]
The policy class \(\Pi\) is a pointwise measurable class of measurable deterministic policies \(\pi:\mathcal X\to\{0,1\}\) with finite VC dimension \(d_\Pi<\infty\).  There exists a countable subclass
\[
\Pi_0=\{\pi_j:j\ge 1\}\subset\Pi
\]
such that every \(\pi\in\Pi\) is the pointwise limit of a sequence of policies from \(\Pi_0\).
\end{assumptionv}

\begin{assumptionv}{ass:optimal-in-class}[Oracle policy inclusion]
For every law \(P\) in the law class, the oracle policy \(\pi^\star_P\) belongs to \(\Pi\).
\end{assumptionv}

\begin{assumptionv}{ass:vc-localized-envelope}[Localized VC envelope]
For the pointwise measurable finite-VC policy class \(\Pi\), every centered policy-indexed process with envelope \(B\) and conditional second moment bounded by
\[
C B^2 P_X(D_\pi)
\]
has fixed-radius localized supremum bounded, under margin localization, by
\[
C B n^{-1/2} r^{\alpha/(2+2\alpha)}(\log n)^p .
\]
\end{assumptionv}

\begin{assumptionv}{ass:margin-window}[Margin window]
The margin window satisfies
\[
0<u_0<2.
\]
\end{assumptionv}

\begin{assumptionv}{ass:nuisance-rate}[Cross-fitted nuisance rates]
The cross-fitted nuisance estimators satisfy, for each \(a\in\{0,1\}\),
\[
\|\widehat\mu_a-\mu_a\|_{L_2(P)}\le r_{\mu,n},
\qquad
\|\widehat e-e\|_{L_2(P)}\le r_{e,n},
\]
and their product rate satisfies
\[
r_{\mu,n}r_{e,n}=O(n^{-1/2}).
\]
\end{assumptionv}

\begin{assumptionv}{ass:strict-overlap-endpoint}[Strict-overlap endpoint]
When \(\gamma=0\), there exists a constant \(\underline p\in(0,1/2]\) such that
\[
p_P(X)\ge \underline p
\]
holds \(P_X\)-almost surely.
\end{assumptionv}

\begin{assumptionv}{ass:bounded-crossfit-nuisances}[Bounded Cross-Fit Nuisances]
The cross-fitted outcome-regression estimators \(\widehat\mu_0\) and \(\widehat\mu_1\) take values in \([-1,1]\).
\end{assumptionv}

\begin{assumptionv}{ass:polynomial-nuisance-exponents}[Polynomial Nuisance Rates]
There exist exponents \(a\ge 0\) and \(c\ge 1/2\), and constants \(C_\mu\) and \(C_{\mathrm{prod}}\), such that, for all sufficiently large \(n\),
\[
r_{\mu,n}\le C_\mu n^{-a},
\qquad
r_{\mu,n}r_{e,n}\le C_{\mathrm{prod}}n^{-c}.
\]
\end{assumptionv}

\begin{assumptionv}{ass:fixed-crossfit-fold-count}[Fixed Cross-Fit Fold Count]
The number \(K\) of cross-fitting folds is a fixed finite integer independent of \(n\), and the deterministic partition \(I_1,\ldots,I_K\) of the observation indices is balanced.
\end{assumptionv}

\begin{assumptionv}{ass:vc-localized-offset-envelope}[Localized Offset Envelope]
For the pointwise measurable finite-VC policy class \(\Pi\), every centered policy-indexed empirical process
\[
z_\pi=(P_n-P)g_\pi
\]
built from an i.i.d. sample of size \(n\), with envelope \(B\) and conditional second moment bounded by
\[
C B^2 P_X(D_\pi),
\]
satisfies
\[
\mathbb E\sup_{\pi\in\Pi}
\left\{2|z_\pi|-\frac{R_P(\pi)}{4}\right\}_+
\le
C\left(\frac{B^2}{n}\right)^{A_\alpha}(\log n)^p .
\]
\end{assumptionv}

\begin{definitionv}{def:welfare-regret}[Welfare \(V_P(\pi)\) and Regret \(R_P(\pi)\)]
For a policy \(\pi\), its welfare under \(P\) is
\[
V_P(\pi)=\mathbb E_P\!\left[\pi(X)\tau_P(X)\right].
\]
The oracle policy is
\[
\pi^\star_P(x)=\mathbf 1\{\tau_P(x)\ge 0\}.
\]
The regret of \(\pi\) under \(P\) is
\[
R_P(\pi)=V_P(\pi^\star_P)-V_P(\pi).
\]
\end{definitionv}

\begin{definitionv}{def:disagreement}[Disagreement Set \(D_\pi\)]
For a policy \(\pi\), its disagreement set is
\[
D_\pi=\{x\in\mathcal X:\pi(x)\ne \pi^\star_P(x)\}.
\]
\end{definitionv}

\begin{definitionv}{def:exponents}[Exponents $\beta_{\alpha,\gamma}$, $D_{\alpha,\gamma}$, and $r_\star(\alpha,\gamma)$]
The overlap calibration exponent, denominator exponent, and minimax exponent are
\[
\beta_{\alpha,\gamma}
=
\begin{cases}
0, & \gamma=0,\\
\dfrac{\alpha\gamma}{\alpha+1}, & \gamma>0,
\end{cases}
\qquad
D_{\alpha,\gamma}=2+\alpha+\beta_{\alpha,\gamma},
\qquad
r_\star(\alpha,\gamma)=\frac{1+\alpha}{D_{\alpha,\gamma}}.
\]
\end{definitionv}

\begin{definitionv}{def:feasible-rate}[Feasible upper exponent $r_{\mathrm{up}}=r_{\mathrm{feas}}$]
Fix one nuisance regime \((a,c,C_\mu,C_{\mathrm{prod}})\), and set
\[
A_\alpha=\frac{1+\alpha}{2+\alpha}.
\]

If \(\gamma>0\), fix constants
\[
\bar u\in(0,u_0],
\qquad
q_0\in\bigl(0,\min\{1/2,c_o\bar u^\gamma\}\bigr].
\]
For \(0\le s\le 1/2\) and \(0\le t\le s/\gamma\), define
\[
\phi(s,t)
=
\min\left\{
A_\alpha(1-2s),\,
c-s,\,
a+\frac{s}{2\gamma}+\frac{\alpha t}{2},\,
2a-t
\right\}.
\]
Let \((s_{\mathrm{feas}},t_{\mathrm{feas}})\) be any maximizer of \(\phi\) over this compact feasible set, and set
\[
g_{\mathrm{joint}}(\alpha,\gamma,a,c)
=
\phi(s_{\mathrm{feas}},t_{\mathrm{feas}}),
\qquad
q_n=q_0 n^{-s_{\mathrm{feas}}},
\qquad
u_n=\bar u n^{-t_{\mathrm{feas}}}.
\]
Then \(q_n\le c_o u_n^\gamma\) for all sufficiently large \(n\), and the solved conditional feasible upper exponent is
\[
r_{\mathrm{up}}
=
r_{\mathrm{feas}}
=
\min\left\{r_\star(\alpha,\gamma),\,g_{\mathrm{joint}}(\alpha,\gamma,a,c)\right\}.
\]

If \(\gamma=0\), take a fixed clipping sequence \(q_n=q_0\) with \(q_0\le \underline p/2\), and set
\[
r_{\mathrm{up}}
=
r_{\mathrm{feas}}
=
\min\{A_\alpha,c\}.
\]
\end{definitionv}

\begin{definitionv}{def:clipped-propensity}[Clipped propensity $e_q(x;\eta)$]
For a nuisance triple \(\eta=(\mu_0,\mu_1,e)\) and a clipping level \(q\in(0,1/2]\), the clipped propensity score is
\[
e_q(x;\eta)
=
\min\{1-q,\max\{q,e(x)\}\}.
\]
\end{definitionv}

\begin{definitionv}{def:clipped-aipw-score}[Clipped AIPW score $\Gamma_q(O;\eta)$]
For \(O=(X,A,Y)\), nuisance triple \(\eta=(\mu_0,\mu_1,e)\), and clipping level \(q\in(0,1/2]\), the clipped augmented inverse-propensity-weighted score for the treatment contrast is
\[
\Gamma_q(O;\eta)
=
\mu_1(X)-\mu_0(X)
+
\frac{A}{e_q(X;\eta)}\{Y-\mu_1(X)\}
-
\frac{1-A}{1-e_q(X;\eta)}\{Y-\mu_0(X)\}.
\]
\end{definitionv}

\begin{definitionv}{def:feasible-erm}[Feasible clipped-AIPW ERM $\widehat\pi_n$]
Fix the countable pointwise dense subclass
\[
\Pi_0=\{\pi_j:j\ge 1\}
\]
of \(\Pi\), and fix a deterministic balanced \(K\)-fold partition
\[
I_1,\ldots,I_K
\]
of \(\{1,\ldots,n\}\), with each fold having size \(\lfloor n/K\rfloor\) or \(\lceil n/K\rceil\). Let \(k(i)\) denote the evaluation fold containing observation \(i\). Given foldwise cross-fitted nuisance estimators \(\widehat\eta^{(-k)}\) and a clipping level \(q\in(0,1/2]\), define the empirical welfare criterion
\[
\widehat V_{n,q}(\pi)
=
\frac1n\sum_{i=1}^n
\pi(X_i)\,
\Gamma_q\!\left(O_i;\widehat\eta^{(-k(i))}\right).
\]
Let \(j_n\) be the smallest index \(j\ge1\) such that
\[
\widehat V_{n,q}(\pi_j)
\ge
\sup_{\ell\ge1}\widehat V_{n,q}(\pi_\ell)-\frac1n,
\]
and set
\[
\widehat\pi_n=\pi_{j_n}.
\]
This measurable \(\Pi\)-valued \(1/n\)-empirical risk maximizer over \(\Pi_0\) is the feasible clipped-AIPW ERM. The clipping level \(q\) is the tuning parameter chosen by the bias-variance balance.
\end{definitionv}

\begin{definitionv}{def:minimax-regret}[Minimax regret $M_n(\alpha,\gamma)$]
The minimax regret over the law class \(\mathcal P_{\alpha,\gamma}\) is
\[
M_n(\alpha,\gamma)
=
\inf_{\widehat\pi_n}
\sup_{P\in\mathcal P_{\alpha,\gamma}}
\mathbb E_P\!\left[
R_P(\widehat\pi_n)
\right],
\]
where the infimum ranges over all measurable data-dependent estimators \(\widehat\pi_n\) taking values in \(\Pi\).
\end{definitionv}

\begin{definitionv}{def:upper-risk}[Conditional upper risk \(U_n(\alpha,\gamma,a,c;\widehat\eta)\)]
Fix a nuisance regime \((a,c,C_\mu,C_{\mathrm{prod}})\) and the associated deterministic schedules \((q_n,u_n)\). The conditional upper risk is
\[
U_n(\alpha,\gamma,a,c;\widehat\eta)
=
\sup E_P\!\left[R_P(\widehat\pi_n)\right],
\]
where the supremum is over all laws \(P\in\mathcal P_{\alpha,\gamma}\) for which the oracle policy belongs to \(\Pi\), all supplied cross-fitted nuisance estimators \(\widehat\eta^{(-k)}\) satisfying the stated \(L_2(P)\) nuisance-rate bounds, boundedness conditions, and polynomial nuisance-rate bounds with exponents \(a,c\) and constants \(C_\mu,C_{\mathrm{prod}}\), all policy classes \(\Pi\) satisfying the stated finite-VC, localized-envelope, and localized-offset-envelope conditions, and all cross-fitting schemes with fixed fold count \(K\).
\end{definitionv}

\begin{definitionv}{def:law-class}[Law class \(\mathcal P_{\alpha,\gamma}\)]
The class \(\mathcal P_{\alpha,\gamma}\) consists of all observed laws \(P\) on \(\Omega\) for which Assumptions~\ref{obj:ass:bounded-outcome}, \ref{obj:ass:positivity}, \ref{obj:ass:margin}, \ref{obj:ass:zero-effect}, \ref{obj:ass:overlap-decay}, and \ref{obj:ass:strict-overlap-endpoint} hold at exponents \((\alpha,\gamma)\).
\end{definitionv}

\begin{definitionv}{def:two-point-witness}[Two-point witness laws \(P_{n,\sigma}\)]
Let \(\mathcal X=[0,1]\) and let \(P_X\) be Lebesgue measure. Set
\[
h_n=n^{-1/(2+\alpha+\beta_{\alpha,\gamma})},
\qquad
q_n=
\begin{cases}
1/4, & \beta_{\alpha,\gamma}=0,\\
h_n^{\beta_{\alpha,\gamma}}, & \beta_{\alpha,\gamma}>0,
\end{cases}
\]
and
\[
B_n=[0,c_Bh_n^\alpha],
\qquad
P_X(B_n)=c_Bh_n^\alpha .
\]
Choose \(\tau_0\in(u_0,2)\). For each \(\sigma\in\{+,-\}\), the law \(P_{n,\sigma}\) has the same covariate law \(P_X\), the same logging propensity
\[
e(x)=
\begin{cases}
q_n, & x\in B_n,\\
1/2, & x\notin B_n,
\end{cases}
\]
and the same off-block outcome law supported on \(\{-1,1\}\), with
\[
E[Y\mid X=x,A=1]=\tau_0/2,
\qquad
E[Y\mid X=x,A=0]=-\tau_0/2,
\qquad x\notin B_n .
\]
Thus the common off-block optimal label is \(1\), and the common off-block contrast is \(\tau_0\). On \(B_n\), set \(Y=0\) almost surely when \(A=0\). On the active treated cell \(\{X\in B_n,A=1\}\), let \(Y\) be supported on \(\{-1,1\}\) with
\[
P(Y=1\mid X\in B_n,A=1)=\frac{1+\sigma h_n}{2},
\qquad
P(Y=-1\mid X\in B_n,A=1)=\frac{1-\sigma h_n}{2}.
\]
Consequently,
\[
E[Y\mid X\in B_n,A=1]=\sigma h_n,
\qquad
\tau_{P_{n,\sigma}}(x)=\sigma h_n
\quad\text{for }x\in B_n .
\]
The two laws agree everywhere except on the treated active cell. The constants satisfy
\[
8C_BC_Q<\log 5 .
\]
\end{definitionv}

\begin{theoremv}{thm:welfare-identity}[Welfare regret identity]
Under Assumption~\ref{obj:ass:bounded-outcome}, for every deterministic policy \(\pi\in\Pi\),
\[
R_P(\pi)
=
E_P\!\left[
|\tau_P(X)|\,\mathbf 1\{\pi(X)\ne\pi^\star_P(X)\}
\right]
=
\int_{\mathcal X}
|\tau_P(x)|\,\mathbf 1\{\pi(x)\ne\pi^\star_P(x)\}\,dP_X(x).
\]
\end{theoremv}

\begin{theoremv}{thm:margin-localization}[Margin localization bound]
Under Assumptions~\ref{obj:ass:margin}, \ref{obj:ass:zero-effect}, and \ref{obj:ass:bounded-outcome}, and with \(D_\pi\) as in Definition~\ref{obj:def:disagreement}, there exists a constant
\[
C=C(C_m,u_0,\alpha)
\]
such that, for every \(\pi\in\Pi\),
\[
P_X(D_\pi)
\le
C\,R_P(\pi)^{\alpha/(1+\alpha)} .
\]
\end{theoremv}

\begin{lemmav}{prop:overlap-envelope}[Overlap envelope calibration]
For \(\alpha\ge0\), \(\gamma>0\), \(h\in(0,1)\), and \(\beta\ge0\), consider the tight window
\[
v=h^\beta,
\qquad
u=h^{\beta/\gamma}.
\]
Then
\[
u^\alpha v^{1/\gamma}
=
h^{(\alpha+1)\beta/\gamma},
\]
and
\[
u^\alpha v^{1/\gamma}\ge h^\alpha
\quad\Longleftrightarrow\quad
\beta\le \frac{\alpha\gamma}{\alpha+1},
\]
with equality at \(\beta=\beta_{\alpha,\gamma}\). Hence a block with mass of order \(h^\alpha\), contrast of order \(h\), and propensity of order \(h^\beta\) satisfies the overlap-decay envelope if and only if
\[
\beta\le\beta_{\alpha,\gamma}.
\]
The exponent \(\beta_{\alpha,\gamma}\) is therefore the least informative admissible weak-arm exponent.
\end{lemmav}

\begin{lemmav}{lem:two-point-divergence}[Two-point divergence]
For the witness laws in Definition~\ref{obj:def:two-point-witness}, under Assumption~\ref{obj:ass:margin-window}, the one-observation chi-square divergence satisfies
\[
\chi^2(P_{n,+}\,\|\,P_{n,-})
 \le C\,m_n q_n h_n^2,
\]
where
\[
m_n=P_X(B_n)\asymp h_n^\alpha,
\qquad
q_n\asymp h_n^{\beta_{\alpha,\gamma}}.
\]
Consequently,
\[
\chi^2(P_{n,+}\,\|\,P_{n,-})
 \lesssim h_n^{2+\alpha+\beta_{\alpha,\gamma}}.
\]
For
\[
h_n=n^{-1/(2+\alpha+\beta_{\alpha,\gamma})},
\]
the product divergence is bounded:
\[
\chi^2(P_{n,+}^{\otimes n}\,\|\,P_{n,-}^{\otimes n})\le C.
\]
Equivalently,
\[
1+\chi^2(P_{n,+}^{\otimes n}\,\|\,P_{n,-}^{\otimes n})
=
\left(1+O\!\left(h_n^{2+\alpha+\beta_{\alpha,\gamma}}\right)\right)^n
\le e^C .
\]
\end{lemmav}

\begin{lemmav}{lem:regret-separation}[Regret separation]
For the witness laws in Definition~\ref{obj:def:two-point-witness}, under Assumption~\ref{obj:ass:margin-window},
\[
\inf_{\pi\in\Pi}\max_{\sigma\in\{+,-\}}
R_{P_{n,\sigma}}(\pi)
\ge
c\,h_n^{1+\alpha}.
\]
The two optimal policies disagree on \(B_n\). Therefore every \(\pi\in\Pi\) misclassifies \(B_n\) under at least one of the two laws, and under that law incurs regret at least
\[
\frac12\,P_X(B_n)\,h_n
\asymp
h_n^{1+\alpha}.
\]
\end{lemmav}

\begin{lemmav}{lem:le-cam-two-point-chisq}[Chi-square testing bound]
If
\[
\chi^2(P_+^{\otimes n}\,\|\,P_-^{\otimes n})\le C_\chi,
\]
then every test \(\psi\) taking values in \(\{+,-\}\) satisfies
\[
P_+^{\otimes n}(\psi\ne +)
+
P_-^{\otimes n}(\psi\ne -)
\ge
c(C_\chi)>0.
\]
In particular, this conclusion holds whenever the one-observation chi-square divergence is bounded by \(c_0/n\), since
\[
1+\chi^2(P_+^{\otimes n}\,\|\,P_-^{\otimes n})
=
\left(1+\chi^2(P_+\,\|\,P_-)\right)^n
\le
\left(1+\frac{c_0}{n}\right)^n
\le e^{c_0}.
\]
\end{lemmav}

\begin{theoremv}{thm:minimax-lower}[Minimax lower bound]
Under Assumption~\ref{obj:ass:margin-window}, for the minimax regret \(M_n(\alpha,\gamma)\) of Definition~\ref{obj:def:minimax-regret} over the law class \(\mathcal P_{\alpha,\gamma}\) of Definition~\ref{obj:def:law-class}, there is a constant \(c>0\), independent of \(n\), such that for all sufficiently large \(n\),
\[
M_n(\alpha,\gamma)
=
\inf_{\widehat\pi}
\sup_{P\in\mathcal P_{\alpha,\gamma}}
\mathbb E_P R_P(\widehat\pi)
\ge
c\,n^{-r_\star(\alpha,\gamma)}.
\]
\end{theoremv}

\begin{theoremv}{oeq:feasible-upper}[Feasible Upper Exponent]
Fix a nuisance regime \((a,c,C_\mu,C_{\mathrm{prod}})\) satisfying Assumption~\ref{obj:ass:polynomial-nuisance-exponents}. Consider the law class \(\mathcal P_{\alpha,\gamma}\) of Definition~\ref{obj:def:law-class}, the feasible cross-fitted clipped-AIPW \(1/n\)-ERM \(\widehat\pi_n\) of Definition~\ref{obj:def:feasible-erm}, and the conditional upper risk \(U_n(\alpha,\gamma,a,c;\widehat\eta)\) of Definition~\ref{obj:def:upper-risk}.

Assume:
\begin{itemize}
\item \textbf{(Policy class.)} The policy class satisfies Assumption~\ref{obj:ass:policy-class}.
\item \textbf{(Cross-fitting.)} The number of folds satisfies Assumption~\ref{obj:ass:fixed-crossfit-fold-count}.
\item \textbf{(Nuisance boundedness.)} The cross-fitted nuisance estimators satisfy Assumption~\ref{obj:ass:bounded-crossfit-nuisances}.
\item \textbf{(Localized envelopes.)} The localized empirical-process and offset envelope conditions in Assumptions~\ref{obj:ass:vc-localized-envelope} and~\ref{obj:ass:vc-localized-offset-envelope} hold.
\end{itemize}

Then the conditional upper risk of \(\widehat\pi_n\), uniformly over \(\mathcal P_{\alpha,\gamma}\) and over the side-condition domain in Definition~\ref{obj:def:upper-risk}, has upper exponent \(r_{\mathrm{up}}\):
\[
U_n(\alpha,\gamma,a,c;\widehat\eta)\le C n^{-r_{\mathrm{up}}}(\log n)^p .
\]
\end{theoremv}

\begin{definitionv}{oeq:feasible-tight}[Feasible Tightness Question]
In the strict-gap branch of the conditional feasible achievability bound for the regime-indexed risk \(U_n\), it is open whether a genuinely feasible estimator using estimated cross-fitted nuisance functions can attain the converse exponent \(r_\star(\alpha,\gamma)\), or whether weak-arm nuisance learning imposes the slower conditional exponent derived by the feasible upper bound.
\end{definitionv}

\begin{lemmav}{lem:feasible-erm-basic-inequality}[Feasible ERM Inequality]
Under Assumption~\ref{obj:ass:policy-class}, Definition~\ref{obj:def:feasible-erm} defines a measurable \(\Pi\)-valued estimator \(\widehat\pi_n\) satisfying
\[
\widehat V_{n,q}(\widehat\pi_n)
\ge
\sup_{\pi\in\Pi}\widehat V_{n,q}(\pi)-\frac1n .
\]
Hence, for every comparator \(\pi^b\in\Pi\),
\[
P_n\!\left[(\widehat\pi_n-\pi^b)\Gamma_q\right]\ge -\frac1n .
\]
Under Assumption~\ref{obj:ass:optimal-in-class}, this conclusion applies with \(\pi^b=\pi^\star_P\).
\end{lemmav}

\begin{lemmav}{lem:crossfit-localized-process-reduction}[Cross-Fit Localized Reduction]
Let \(Z_{\mathrm{cf}}(r)\) denote the pooled cross-fit localized supremum formed by averaging centered evaluation-fold increments.

Assume:
\begin{itemize}
\item \textbf{(Data and folds.)} The observations are i.i.d. as in Assumption~\ref{obj:ass:iid}, the fixed balanced cross-fitting scheme of Assumption~\ref{obj:ass:fixed-crossfit-fold-count} is used, and the risk is evaluated in the setting of Definition~\ref{obj:def:upper-risk}.
\item \textbf{(Policies and localization.)} The policy class satisfies Assumption~\ref{obj:ass:policy-class}, the margin and zero-effect conditions of Assumptions~\ref{obj:ass:margin} and~\ref{obj:ass:zero-effect} hold, and disagreements are measured by \(D_\pi\) as in Definition~\ref{obj:def:disagreement}.
\item \textbf{(Outcome and envelope.)} The bounded-outcome condition of Assumption~\ref{obj:ass:bounded-outcome} and the localized VC envelope condition of Assumption~\ref{obj:ass:vc-localized-envelope} hold.
\item \textbf{(Foldwise increments.)} Conditional on the training folds, the centered foldwise increments have envelope \(B\) and conditional second moment bounded by
\[
C B^2 P_X(D_\pi).
\]
\end{itemize}
Then, conditional on the training folds, each balanced evaluation fold is i.i.d., and for every \(r\),
\[
E_P\!\left[ Z_{\mathrm{cf}}(r)\mid \text{training folds}\right]
\le
C B n^{-1/2} r^{\alpha/(2+2\alpha)}(\log n)^p .
\]
Consequently,
\[
E_P Z_{\mathrm{cf}}(r)
\le
C B n^{-1/2} r^{\alpha/(2+2\alpha)}(\log n)^p .
\]
\end{lemmav}

\begin{lemmav}{lem:crossfit-localized-offset-control}[Cross-Fit Offset Control]
Let \(G_{\mathrm{cf}}(\pi)\) be the pooled cross-fitted centered policy-indexed process. Suppose that, conditionally on the nuisance fits, the foldwise centered increments have envelope \(B\) and conditional second moment bounded by
\[
C B^2 P_X(D_\pi),
\]
where \(D_\pi\) is the disagreement set of Definition~\ref{obj:def:disagreement}.

Assume:
\begin{itemize}
\item \textbf{(Sampling.)} The observations are i.i.d. as in Assumption~\ref{obj:ass:iid}.
\item \textbf{(Bounded outcomes.)} The outcome is bounded as in Assumption~\ref{obj:ass:bounded-outcome}.
\item \textbf{(Policy class.)} The policy class satisfies Assumption~\ref{obj:ass:policy-class}.
\item \textbf{(Cross-fitting.)} The number of folds satisfies Assumption~\ref{obj:ass:fixed-crossfit-fold-count}.
\item \textbf{(Margin and zero effect.)} The margin and zero-effect conditions in Assumptions~\ref{obj:ass:margin} and~\ref{obj:ass:zero-effect} hold.
\item \textbf{(Offset envelope.)} The localized offset envelope condition in Assumption~\ref{obj:ass:vc-localized-offset-envelope} holds.
\end{itemize}

Then
\[
\mathbb E_P
\left[
\sup_{\pi\in\Pi}
\left\{
2|G_{\mathrm{cf}}(\pi)|-\frac{R_P(\pi)}{4}
\right\}_{+}
\right]
\le
C\left(\frac{B^2}{n}\right)^{A_\alpha}(\log n)^p .
\]
\end{lemmav}

\begin{lemmav}{lem:crude-clipped-score-envelope}[Clipped-Score Envelope]
Under the bounded-outcome condition of Assumption~\ref{obj:ass:bounded-outcome}, for every \(q\in(0,1/2]\) and every cross-fitted nuisance triple \(\widehat\eta^{(-k)}\), the clipped AIPW score of Definition~\ref{obj:def:clipped-aipw-score} satisfies
\[
\left|\Gamma_q\!\left(O;\widehat\eta^{(-k)}\right)\right|\le \frac{C}{q}
\quad\text{almost surely}.
\]
Consequently,
\[
E\!\left[\Gamma_q\!\left(O;\widehat\eta^{(-k)}\right)^2\mid X\right]
\le \frac{C}{q^2}.
\]
\end{lemmav}

\begin{lemmav}{lem:clipped-region-localization}[Clipped-Region Localization]
Assume \(\gamma>0\). Under the overlap-decay and zero-effect conditions of Assumptions~\ref{obj:ass:overlap-decay} and~\ref{obj:ass:zero-effect}, together with the bounded-outcome condition of Assumption~\ref{obj:ass:bounded-outcome}, every policy \(\pi\) with regret \(r=R_P(\pi)\) satisfies, for its disagreement set \(D_\pi\) from Definition~\ref{obj:def:disagreement},
\[
P_X\!\left(D_\pi\cap\{x:p_P(x)\le q\}\right)
\le
C u^\alpha q^{1/\gamma}+\frac{r}{u}
\]
whenever \(0<u\le u_0\) and \(q\le c_o u^\gamma\).
\end{lemmav}

\begin{lemmav}{lem:clip-balance-exponent}[Clipping Balance Exponent]
Under the polynomial nuisance-rate exponents of Assumption~\ref{obj:ass:polynomial-nuisance-exponents}, the deterministic clipping and localization schedules satisfy the following exponent balance.

For \(\gamma>0\), let
\[
A_\alpha=\frac{1+\alpha}{2+\alpha},
\qquad
q_n=q_0 n^{-s_{\mathrm{feas}}},
\qquad
u_n=\bar u\, n^{-t_{\mathrm{feas}}},
\]
and
\[
g_{\mathrm{joint}}
=
\max_{\substack{0\le s\le 1/2\\ 0\le t\le s/\gamma}}
\min\left\{
A_\alpha(1-2s),\,
c-s,\,
a+\frac{s}{2\gamma}+\frac{\alpha t}{2},\,
2a-t
\right\}.
\]
Then the master bound is at most
\[
C\left\{n^{-r_\star(\alpha,\gamma)}+n^{-g_{\mathrm{joint}}}\right\}(\log n)^p
=
C n^{-\min\{r_\star(\alpha,\gamma),g_{\mathrm{joint}}\}}(\log n)^p .
\]
Equivalently, the optimized upper-bound exponent is
\[
r_{\mathrm{up}}=r_{\mathrm{feas}}
=
\min\{r_\star(\alpha,\gamma),g_{\mathrm{joint}}\}.
\]

For \(\gamma=0\), the corresponding exponent is
\[
\min\{A_\alpha,c\}.
\]
\end{lemmav}

\begin{lemmav}{lem:crude-localized-master-bound}[Crude Localized Master Bound]
Let \(\widehat\pi_n\) be the cross-fitted clipped-AIPW \(1/n\)-ERM of Definition~\ref{obj:def:feasible-erm}. Suppose \(P\in\mathcal P_{\alpha,\gamma}\), where \(\mathcal P_{\alpha,\gamma}\) is the law class of Definition~\ref{obj:def:law-class}.

Assume:
\begin{itemize}
\item \textbf{(Policy and oracle.)} The policy class and oracle satisfy Assumptions~\ref{obj:ass:policy-class} and~\ref{obj:ass:optimal-in-class}; welfare and regret are as in Definition~\ref{obj:def:welfare-regret}.
\item \textbf{(Sampling.)} The observations are i.i.d. as in Assumption~\ref{obj:ass:iid}.
\item \textbf{(Nuisance rates.)} The nuisance estimators satisfy Assumption~\ref{obj:ass:nuisance-rate} and the polynomial nuisance exponent condition in Assumption~\ref{obj:ass:polynomial-nuisance-exponents}.
\item \textbf{(Cross-fitting and boundedness.)} The number of folds is fixed as in Assumption~\ref{obj:ass:fixed-crossfit-fold-count}, and the cross-fitted nuisance estimators satisfy Assumption~\ref{obj:ass:bounded-crossfit-nuisances}.
\item \textbf{(Localized empirical processes.)} The localized envelope and offset envelope conditions in Assumptions~\ref{obj:ass:vc-localized-envelope} and~\ref{obj:ass:vc-localized-offset-envelope} hold.
\item \textbf{(Basic inequality, bias, and process bounds.)} The feasible ERM basic inequality of Lemma~\ref{obj:lem:feasible-erm-basic-inequality}, the clipped-bias identity of Lemma~\ref{obj:lem:clip-bias}, the cross-fit offset bound of Lemma~\ref{obj:lem:crossfit-localized-offset-control}, and the localized VC self-bound of Lemma~\ref{obj:lem:localized-vc-self-bound} hold.
\end{itemize}

If \(\gamma>0\) and
\[
q\le c_o u^\gamma,
\]
then
\[
\mathbb E_P R_P(\widehat\pi_n)
\le
C\left\{
n^{-r_\star}
+(nq^2)^{-A_\alpha}
+\frac{r_{\mu,n}r_{e,n}}{q}
+r_{\mu,n}u^{\alpha/2}q^{1/(2\gamma)}
+\frac{r_{\mu,n}^2}{u}
\right\}(\log n)^p .
\]

If \(\gamma=0\) and \(q\le \underline p/2\) is fixed, then
\[
\mathbb E_P R_P(\widehat\pi_n)
\le
C\left\{
n^{-A_\alpha}+r_{\mu,n}r_{e,n}
\right\}(\log n)^p .
\]
\end{lemmav}

\begin{assumptionv}{ass:overlap-decay}[Overlap decay]

For all \(u\) and \(v\) with \(0<u\le u_0\) and \(v>0\) satisfying
\[
v\le c_o u^\gamma,
\]
the joint small-contrast and weak-overlap probability obeys
\[
P\{p_P(X)\le v,\ 0<|\tau_P(X)|\le u\}
   \le C_o u^\alpha v^{1/\gamma},
\]
with the convention that \(v^{1/\gamma}=1\) when \(\gamma=0\).

\end{assumptionv}
\begin{lemmav}{lem:witness-membership}[Witness-law membership]

Assume \(\alpha\ge0\), \(\gamma\ge0\), the margin-window normalization \(0<u_0<2\) in Assumption~\ref{obj:ass:margin-window}, \(C_m,C_o,c_o,c_B,\underline p>0\), \(c_B\le C_m\), \(c_B\le C_o\), and \(\underline p\le 1/4\). Also assume the witness constant satisfies the overlap-decay smallness conditions
\[
0<\gamma,\ 0<\alpha
\quad\Longrightarrow\quad
c_B\le C_o c_o^{-\alpha/\gamma},
\qquad
0<\gamma,\ \alpha=0
\quad\Longrightarrow\quad
c_B\le C_o 4^{-1/\gamma}.
\]
Then, for each witness sign \(\sigma\in\{+,-\}\), for all sufficiently large \(n\), the witness law \(P_{n,\sigma}\) of Definition~\ref{obj:def:two-point-witness} satisfies the observed-law membership conditions defining the class \(\mathcal P_{\alpha,\gamma}\) in Definition~\ref{obj:def:law-class}. Hence, for all sufficiently large \(n\),
\[
P_{n,+},P_{n,-}\in \mathcal P_{\alpha,\gamma}.
\]
Moreover, the two corresponding witness-optimal policies are
\[
x\mapsto 1
\qquad\text{and}\qquad
x\mapsto \mathbf 1\{x\notin B_n\},
\]
where \(B_n=\{x:0\le x\le c_B h_n^\alpha\}\). If these two policies belong to \(\Pi\), then they are available as the two policy actions in the associated two-point lower-bound reduction. The construction uses the off-block contrast \(\tau_0=(u_0+2)/2\), so \(\tau_0\in(u_0,2)\) under the normalization \(0<u_0<2\).

\end{lemmav}
\begin{theoremv}{thm:rate-characterization}[Minimax Lower Rate]

Under Assumption~\ref{obj:ass:margin-window}, let \(\alpha,\gamma\ge0\) and suppose the constants satisfy \(C_m,C_o,c_o,c_B,\underline p>0\), \(c_B\le C_m\), \(c_B\le C_o\), \(\underline p\le 1/4\), and \(8c_B<\log 5\). In addition, assume the overlap-calibration inequalities
\[
0<\gamma,\ 0<\alpha \implies c_B\le C_o c_o^{-\alpha/\gamma},
\qquad
0<\gamma,\ \alpha=0 \implies c_B\le C_o 4^{-1/\gamma}.
\]
If \(\Pi\) is nonempty and every \(\pi\in\Pi\) is measurable, then for the minimax regret over the law class in Definitions~\ref{obj:def:minimax-regret} and~\ref{obj:def:law-class}, there exists a constant \(c>0\) such that, for all sufficiently large \(n\),
\[
M_n(\alpha,\gamma)\ge c\,n^{-r_\star(\alpha,\gamma)},
\qquad
r_\star(\alpha,\gamma)
=
\frac{1+\alpha}{2+\alpha+\beta_{\alpha,\gamma}} .
\]

\end{theoremv}
\begin{lemmav}{lem:clip-bias}[Clipped Score Drift]

Under Assumptions~\ref{obj:ass:bounded-outcome}, \ref{obj:ass:positivity}, and~\ref{obj:ass:bounded-crossfit-nuisances}, and with the clipped propensity and clipped AIPW score of Definitions~\ref{obj:def:clipped-propensity} and~\ref{obj:def:clipped-aipw-score}, let \(\bar\eta=(\bar\mu_0,\bar\mu_1,\bar e)\) be any measurable plug-in nuisance triple. Define
\[
\bar e_q(x)=\min\{1-q,\max\{q,\bar e(x)\}\},
\qquad
\Delta_a(x)=\bar\mu_a(x)-\mu_a(x),
\qquad
b_q(x)=(\bar e_q(x)-e_P(x))\left\{\frac{\Delta_1(x)}{\bar e_q(x)}+\frac{\Delta_0(x)}{1-\bar e_q(x)}\right\}.
\]
Then the clipped score reproduces the contrast up to the drift \(b_q\) in the tested (population) sense: for every bounded measurable \(\varphi\),
\[
\mathbb E_P[\varphi(X)\,\Gamma_q(O;\bar\eta)]-\mathbb E_P[\varphi(X)\,\tau_P(X)]
=
\mathbb E_P[\varphi(X)\,b_q(X)],
\]
equivalently \(\mathbb E_P[\Gamma_q(O;\bar\eta)\mid X=x]-\tau_P(x)=b_q(x)\) for \(P_X\)-almost every \(x\).
In particular, the drift vanishes whenever \(\bar e_q(x)=e_P(x)\) or \(\Delta_0(x)=\Delta_1(x)=0\), since either makes a factor of \(b_q(x)\) vanish; these conditions are sufficient but not necessary, as the bracketed term \(\Delta_1(x)/\bar e_q(x)+\Delta_0(x)/(1-\bar e_q(x))\) may itself vanish. Cancellation does not, however, follow merely from \(p_P(x)>q\): there exist laws and covariate values with \(p_P(x)>q\) at which the drift is nonzero.

\end{lemmav}
\begin{lemmav}{lem:localized-vc-self-bound}[Localized Self Bound]

Let
\[
\rho_n=\left(\frac{B^2}{n}\right)^{A_\alpha}(\log n)^p,
\qquad A_\alpha=\frac{1+\alpha}{2+\alpha},
\qquad B\ge 1,
\qquad p\ge 0,
\]
and let \(C>0\) be the constant of the cross-fit localized offset bound, so that the offset positive part of the centered \(\Pi\)-indexed process \(z_\pi\) is controlled at scale \(C\rho_n\). Assume the sample size is large enough that
\[
\frac1n\le C\rho_n .
\]
Then every \(\Pi\)-indexed estimator \(\widetilde\pi\) satisfying
\[
R_P(\widetilde\pi)\le 2|z_{\widetilde\pi}|+\delta_n,
\qquad \delta_n\ge 0,
\]
also satisfies
\[
\mathbb E_P R_P(\widetilde\pi)\le \tfrac{8}{3}\,(C\rho_n+\delta_n).
\]
If, in addition, \(\delta_n\le 1/n\), then---using \(1/n\le C\rho_n\)---the \(\delta_n\) term is absorbed into the leading rate and
\[
\mathbb E_P R_P(\widetilde\pi)\le \tfrac{8}{3}\,C\,\rho_n .
\]

\end{lemmav}
\begin{lemmav}{lem:localized-vc-process-bound}[Localized VC Process]

Under Assumptions~\ref{obj:ass:policy-class}, \ref{obj:ass:margin}, \ref{obj:ass:zero-effect}, \ref{obj:ass:bounded-outcome}, and~\ref{obj:ass:vc-localized-envelope}, and using the regret risk and disagreement sets in Definition~\ref{obj:def:upper-risk} and Definition~\ref{obj:def:disagreement}, let
\[
\kappa=\frac{\alpha}{2+2\alpha}.
\]
There exist constants \(C>0\) and \(p\ge0\) such that the following holds. For an i.i.d.\ sample of size \(m\ge1\), any envelope \(B\ge0\), and any policy-compatible increment family \(g_\pi\)---that is, \(g_\pi(O)\) depends on the policy \(\pi\) only through its decision \(\pi(X)\)---satisfying
\[
|g_\pi(O)|\le B\quad(\forall\,\pi\in\Pi,\ \forall\,O),
\qquad
\int g_\pi(O)^2\,dP(O)\le B^2\,P_X(D_\pi)\quad(\forall\,\pi\in\Pi),
\]
define
\[
Z_m(r)
=
\sup_{\pi\in\Pi:\,R_P(\pi)\le r}
\left|(P_m-P)g_\pi\right|.
\]
Then, for every regret radius \(r\ge0\),
\[
\mathbb E_P Z_m(r)
\le
C B\,m^{-1/2} r^\kappa(\log m)^p .
\]

\end{lemmav}
\begin{lemmav}{lem:localized-clipped-drift-bound}[Localized Clipped Drift Bound]

Under Assumptions~\ref{obj:ass:nuisance-rate}, \ref{obj:ass:polynomial-nuisance-exponents}, \ref{obj:ass:overlap-decay}, \ref{obj:ass:zero-effect}, \ref{obj:ass:bounded-outcome}, and \ref{obj:ass:strict-overlap-endpoint}, and Lemma~\ref{obj:lem:clipped-region-localization}, let \(P\) be a well-formed observed law, let every \(\pi\in\Pi\) be measurable, and fix a clipping level \(q\) with \(0<q\le 1/2\).  With the cross-fitted nuisance estimators fixed on the evaluation fold, write \(\widehat\eta^{(-k)}=(\widehat\mu_0^{(-k)},\widehat\mu_1^{(-k)},\widehat e^{(-k)})\), and suppose \(r_{\mu,n},r_{e,n}\ge0\),
\[
\int (\widehat\mu_0^{(-k)}(x)-\mu_0(x))^2\,dP_X(x)\le r_{\mu,n}^2,
\quad
\int (\widehat\mu_1^{(-k)}(x)-\mu_1(x))^2\,dP_X(x)\le r_{\mu,n}^2,
\]
\[
\int (\widehat e^{(-k)}(x)-e_P(x))^2\,dP_X(x)\le r_{e,n}^2,
\]
and that these three nuisance-error functions belong to \(L_2(P_X)\).  Define
\[
b_q(x)
=
\mathbb E_P\!\left[\Gamma_q(O;\widehat\eta^{(-k)})\mid X=x\right]
-
\tau_P(x).
\]

If \(\gamma>0\), set \(C_0=4\{1+\sqrt{\max\{C_o,1\}}\}\).  Then, for every \(\pi\in\Pi\) and every localization window \(u\) satisfying \(0<u\le u_0\), writing \(r=R_P(\pi)\),
\[
\left|P\!\left[(\pi-\pi^\star_P)b_q\right]\right|
\le
C_0\left\{
\frac{r_{\mu,n}r_{e,n}}{q}
+
r_{\mu,n}u^{\alpha/2}q^{1/(2\gamma)}
+
r_{\mu,n}\left(\frac{r}{u}\right)^{1/2}
\right\}
\]
whenever \(q\le c_o u^\gamma\).

If \(\gamma=0\) and \(q\le \underline p/2\), set \(C_1=\max\{1,2/q\}\).  Then, for every \(\pi\in\Pi\),
\[
\left|P\!\left[(\pi-\pi^\star_P)b_q\right]\right|
\le
C_1 r_{\mu,n}r_{e,n}.
\]

\end{lemmav}

